\begin{abstract}
Long-video anomaly understanding requires broad scene context and sensitivity to brief, sparse event evidence, yet a single uniformly sampled view can underemphasize the moments that determine the final decision.
We introduce \textit{A}nomaly-\textit{T}argeted \textit{R}etrieval (\textit{ATR}), a training-free selective re-reading framework.
Scout first searches the complete video for candidate evidence; Sniper then jointly receives the same complete video and the retrieved clip, preserving global context while adding a denser second view without parameter updates or target-task temporal supervision.
On the official UCF-Crime test set, ATR raises a 4B model to 85.17\% accuracy and reduces its false-negative rate from 26.43\% to 11.43\%.
An 8B configuration reaches 85.52\% accuracy with 9.33\% false-positive rate, outperforming a 32B global-only baseline at about one tenth of its GPU-hour cost.
Compute- and content-matched controls show that a second global call alone provides no gain, while Scout significantly outperforms duration-matched Random and Uniform clips in the 4B setting.
Removing the retrieved view reverses all analyzed rescued decisions, confirming that the final prediction uses the selected evidence.
Together, these results establish content-aware global--local retrieval as a practical and efficient inference strategy for long-video anomaly understanding.
\end{abstract}
    
